%% Separately authored historical derivative for the distinct ccc/name-counting results.
%% Necessity: exact source lines 10675--11353 restore the core forcing proofs
%% inline, while the later ccc, nice-name, and many-Cohen-real proofs remain distinct.
%% Definition: this file is based on historical annex slice 90--178 with an
%% explicit bridge to the renewed body forcing convention.
%% Standing: this is a reciprocal editorial proof annex, not inherited source.
%% Nonredundancy: all overlapping annex items are retired; the retained slice
%% proves only the separately routed ccc/name-counting body results.
%% Use: annex/p06-route.tex inputs this file once and its three identities are
%% reciprocal targets of the bounded body derivative.

% BEGIN HISTORICAL SLICE 90-178
\noindent\textit{Proof standing.}  In all three items below, \(M\),
\(\mathbb P\), \(G\), ground-model cardinal computations, semantic forcing,
and the imported recursive forcing assertion have exactly the meanings fixed
in Convention~\ref{def:p06-ccc-forcing-standing}.  The proofs use the cited
forcing theorem's definability and truth clauses; they do not define a third
forcing relation.

\begin{annexitem}{annex:p6.force.9}{thm:p06-ccc-preservation}{Cardinal and cofinality preservation by ccc forcing}
\begin{proof}
Work in \(M\), using Choice.  First note that if \(D\in M\) is dense below
some \(d\in G\), then
\[
  D\cup\{q\in\mathbb P:q\text{ is incompatible with }d\}
\]
is a globally dense member of \(M\).  Genericity and \(d\in G\) therefore
give \(G\cap D\neq\varnothing\).

If \(M[G]\) collapsed an uncountable
\(M\)-cardinal \(\kappa\), there would be an ordinal
\(\lambda<\kappa\), a surjection \(f:\lambda\to\kappa\) in \(M[G]\), a
name \(\nu_f\in M\) with \(\nu_f^G=f\), and \(p\in G\) such that
\(p\forces \nu_f:\check\lambda\twoheadrightarrow\check\kappa\).  For each
\(\alpha<\lambda\), choose below \(p\) a maximal antichain deciding
\(\nu_f(\check\alpha)\).  The ccc makes it countable; if \(X_\alpha\) is its set
of possible values, the downward closure of that antichain is dense below
\(p\).  The observation above shows that the interpreted map satisfies
\[
  \operatorname{ran}(f)\subseteq
  \bigcup_{\alpha<\lambda}X_\alpha.
\]
In \(M\), the union has size at
most \(\max((\abs\lambda)^M,\aleph_0)<\kappa\), impossible.  The finite cardinals
and \(\omega\) are absolute because the extension has the same ordinals.

Now let \(\mu=\operatorname{cf}^{M}(\delta)\).  A ground-model cofinal map
of type \(\mu\) remains cofinal because the ordinals and their order are
unchanged, so cofinality cannot increase.  If a cofinal map in \(M[G]\) had
infinite domain ordinal \(\lambda<\mu\), choose a name and a condition in
\(G\) forcing this, and repeat the deciding-antichain construction below that
condition.  Since \(\mu\) is regular, the resulting union of possible values
has size \(\max((\abs\lambda)^M,\aleph_0)<\mu\) in \(M\), hence is bounded in
\(\delta\), contradicting cofinality.  A finite range is bounded in every
limit ordinal, so ground cofinality \(\omega\) cannot drop to a finite value;
successor and zero cofinalities are fixed by the unchanged ordinals.
\end{proof}
\end{annexitem}

\begin{annexitem}{annex:p6.force.10}{thm:p06-nice-name-counting}{Nice-name reduction and counting}
\begin{proof}
Fix \(p\forces\nu\subseteq\check\omega\).  For each \(n<\omega\), choose
below \(p\) a maximal antichain \(A_n\) among conditions forcing
\(\check n\in\nu\).  Put
\[
  \zeta=\bigcup_{n<\omega}
  \{\langle\check n,a\rangle:a\in A_n\}.
\]
For any generic \(H\ni p\), a member of \(A_n\cap H\) forces
\(n\in\zeta^H\cap\nu^H\).  Conversely, if \(n\in\nu^H\), the truth lemma
gives \(q\in H\) forcing \(\check n\in\nu\).  Directedness supplies
\(d\in H\) with \(d\leq p,q\), which still forces membership.  Maximality
makes the common extensions of \(d\) with members of \(A_n\) dense below
\(d\).  Adjoining every condition incompatible with \(d\) makes this set a
globally dense member of \(M\).  Since \(d\in H\), genericity cannot meet the
incompatible part, so it supplies a common extension in \(H\); upward closure
then puts the corresponding member of \(A_n\) in \(H\).  Hence
\(n\in\zeta^H\).  Thus
\(p\) semantically forces \(\nu=\zeta\), and the forcing theorem yields
\(p\forces\nu=\zeta\).

Under ccc each \(A_n\) is countable.  There are at most
\(\lambda=(\abs{\mathbb P}^{\aleph_0})^M\) countable subsets of
\(\mathbb P\) in \(M\), and a countable sequence of them has the same bound.
For any \(x\subseteq\omega\) in \(M[G]\), choose \(\nu\in M\) with
\(x=\nu^G\).  The truth lemma supplies \(p\in G\) forcing
\(\nu\subseteq\check\omega\), so the local result gives a nice \(\zeta\)
with \(x=\zeta^G\).  Thus the nice names cover all reals; ccc preservation
keeps \(\lambda\) a cardinal, proving the bound.
\end{proof}
\end{annexitem}

\begin{annexitem}{annex:p6.force.11}{thm:p06-many-cohen-continuum}{Forcing the size of the continuum}
\begin{proof}
For \(\alpha<\kappa\), let
\(c_\alpha(n)=(\bigcup G)(\alpha,n)\).  Coordinate-density makes each
\(c_\alpha\) total.  For distinct \(\alpha,\beta\), conditions assigning
opposite values at a fresh common coordinate are dense, so
\(c_\alpha\neq c_\beta\); hence the
extension has at least \(\kappa\) reals.  The poset has cardinality \(\kappa\),
is ccc by Theorem~\ref{thm:p06-add-ccc}, and preserves \(\kappa\) by
Theorem~\ref{thm:p06-ccc-preservation}.  Nice-name counting and
\(\kappa^{\aleph_0}=\kappa\) give at most \(\kappa\) reals.  Therefore the
continuum is \(\kappa\); for ground \(\mathsf{GCH}\) and
\(\kappa=\aleph_2\), CH fails.
\end{proof}
\end{annexitem}

% END HISTORICAL SLICE 90-178
